% critique.tex — f164r, vv. 123–164 : commentaire critique (articles textuels du commentaire de l'auteur, repris mot pour mot ; voir repartition.md)
\cc{134}{νεανίσχ᾽ εἴ τις εἶ.}{Le manuscrit de Ravenne comporte une coquille : le κ de νεανίσκ᾽ est aspiré alors qu'il est suivi de εἴ. Plusieurs solutions ont été apportées : lire ἥτις (justifiant le féminin par Hésychius (η 948) glosant ἥτις· ὅστις). D'autres se sont contentés de substituer un κ au χ. J'ai une préférence pour cette dernière correction\nc{Signalons qu'une coquille analogue (οὐχ᾽ εἴρηχ᾽ R en place d'un οὐκ εἴρηχ᾽) se retrouve au vers 556. De plus, les anciens manuscrits sur lesquels on a copié les \italique{Thesmophories} n'étaient pas systématiquement élidés ; ce genre de fautes s'explique donc aisément. Cf. par exemple v. 148 ἐσθητ' ἅμα ; v. 279 θραττ' έπου ; v. 1178 ἔρχετ' ὡς.} qui me semble être une pique contre Agathon : « jeune homme... si tu en es un ! » (voir le commentaire de la pièce).}
\cc{139–44}{Fragments d'un ancien papyrus.}{L'on possède quelques fragments d'un ancien papyrus (\PsiUn \ dans mon apparat) du premier ou second siècle ap. J.-C. pour les vers 133-144 ; 145-156 ; 237-245 ; 273-288 ; 594-596 ; 804-809. Il nous apprend en outre le \latin{numerus versuum} d'une ancienne édition des \italique{Thesmophories} : la pièce comportait alors 1201 vers (contre 1231 pour celle de Brunck sur laquelle se basent tous les éditeurs modernes). Le papyrus est reproduit en annexe. Cf. \italique{Papiri della Società Italiana} XI, n° 1194 et  \italique{Papiri della Società Italiana} XIV, p. XV.}
\cc{140}{κατόπτρου.}{La variante κάτροπτον (avec métathèse du ρ) est la seule attestée par l'épigraphie pour les inscriptions contemporaines de la représentation des \italique{Thesmophories}\nc{Voir par exemple les \italique{Inscriptiones Graecae} XXII.1471.47, 1544.58.
Le lecteur pourra aisément vérifier ces dires par lui-même en papillonnant sur les \italique{Inscriptiones Graecae} à l'adresse suivante : \url{https://epigraphy.packhum.org/}}. Ce n'est qu'à partir de la fin du \nieme{IV} siècle que la forme κάτοπτρον, plus conforme à l'étymologie, semble prendre le dessus.}
\cc{161–3}{Ἴβυκος ἐκεῖνος κἀνακρέων ὁ Τήιος / κἀλκαῖος, οἳ περὶ ἁρμονίαν ἐχύμησαν / ἐμιτροφόρουν τε καὶ διεκίνων Ἰωνικῶς,}{Le commentaire ancien le plus pertinent sur notre passage n'est pas écrit en marge du \italique{Ravennas} -- comme c'est hélas bien souvent le cas pour les \italique{Thesmophories} -- mais se trouve dans la \italique{Souda} (ε 989 -- et quasiment à l'identique dans le lexique de Zonaras, p. 707) :\citComm{ἀντὶ τοῦ μίτραις\nc{Zonaras -- ou le \italique{pseudo}-Zonaras pour les esprits chagrins -- comporte un accusatif en place de l'ablatif : ἀντὶ τοῦ μίτρας ἤτοι ζώνας διαζωσάμενοι.} τουτέστι ζώναις διαζωσάμενοι ἁβρῶς περιεπάτουν.\\
Signifie : ceints de mitres -- c'est-à-dire de ceintures -- ils vivaient de façon luxurieuse.}Ce commentaire est doublement intéressant : d'abord, il précise que la μίτρα dont il est question est à prendre au sens de \italique{ceinture} et non de bandeau. Ensuite parce qu'il glose trois des mots du vers 163 : ἐμιτροφόρουν par μίτραις διαζωσάμενοι ; Ἰωνικῶς par ἁβρῶς\nc{Voir Érasme, \italique{Adages} (n° 469, \latin{Risus Ionicus}).} ; et -- plus intéressant -- le \italique{contra metrum} †διεκίνων† (†διεκίνουν† dans la \italique{Souda}) par περιεπάτουν (ici avec le sens figuré de \italique{se conduire}, \italique{vivre de telle} ou \italique{telle manière}). 
Cette glose est donc incompatible avec l'émendation de Toup, très largement reprise dans les éditions modernes. Toup proposa la correction suivante : διεκλῶντ' (de διακλάω, \italique{briser en deux}, d'où \italique{briser une vie par la mollesse} (?!)).
Il est bien préférable -- tant du point de vue paléographique que du sens -- de lire διέκειντ', de διάκειμαι, \italique{être disposé}, synonyme de περιπατέω, et qui se construit fréquemment avec un adverbe (ici : διάκειμαι Ἰωνικῶς, \italique{avoir des dispositions à la vie ionienne}, c'est-à-dire à la mollesse et la luxure).\\
Les sous-entendus d'Agathon sont à peine voilés : les trois poètes pédérastes portent la ceinture à la mode ionienne, c'est-à-dire lâche, voire franchement dénouée...
La métaphore est la suivante : tels qu'ils portaient la ceinture, béante, telles étaient leurs \italique{harmonies}. Le vers 162, οἳ περὶ ἁρμονίαν †ἐχύμησαν†\nc{La rêverie du scholiaste reprise par Suidas, ἐχύμισαν· ἀντὶ τοῦ ἔγχυμον ἐποίησαν, n'est évidemment pas à prendre au sérieux. Pas plus que l'invention du verbe χυμίζω... à moins qu'on ne m'explique à quelle sauce (χυμός) tourner le passage pour qu'il ait un rapport avec les inflexions de voix (\italique{harmonies}). Ces divagations auront probablement troublé davantage encore les scribes, déjà en butte à un texte corrompu.}, se laisse maintenant comprendre. Il faut lire :\citComm{οἳ περὶ ἁρμονίαν ἐκεχήνεσαν}La correction est très légère : une nouvelle faute de iotacisme (entre η et υ), de confusion μ/ν (quasiment identiques sur notre \italique{Ravennas}, comme sur beaucoup de manuscrits écrits en minuscules) et d'une haplographie bien explicable :  $\text{εκεχ} \xrightarrow{\text{saut du même au m. }} \text{ εχ}$.\\
χάσκειν signifie littéralement \italique{avoir la bouche grande ouverte}\nc{Bien souvent : \italique{être béat d'admiration}. Voir Taillardat §§472 et 578.}. En bâillant de la sorte aux corneilles, l'on pouvait émettre de légers sons, comparables aux coquilles grillant sur des chardons\nc{ἀνέχασκον εἷς ἕκαστος ἐμφερέστατα ὀπτωμέναις κόγχαισιν ἐπὶ τῶν ἀνθράκων. Aristophane, \italique{fr.} 67 K.-A.} ou à cette vieille lyre de Cratinos, le rival hors-d'âge d'Aristophane, chevrotant ses derniers et mauvais accords (\italique{Cavaliers}, v. 533)\nc{C'est bien ainsi qu'il faut interpréter τῶν θ' ἁρμονιῶν διαχασκουσῶν : « émettre des sons, ou accords » et non « aux joints béants » (en parlant de la lyre), comme on peut le lire.} :\citComm{τῶν θ' ἁρμονιῶν διαχασκουσῶν.}(Noter la collocation entre (δια-)χάσκω et ἁρμονία.)  C'est la même expression qu'on retrouve dans la bouche d'Agathon : περὶ ἁρμονίαν ἐκεχήνεσαν, « béer une harmonie », c'est-à-dire bêler -- ou chevroter, telle est l'image en français, point très éloignée du grec -- quelques paroles en musique\nc{ἁρμονία a le sens d'arrangement de paroles pour un morceau de musique dans Platon, \italique{Theæt.} 175e-176a.}. Le sel de la métaphore réside dans le double sens de περιχάσκειν. Le verbe est justement attesté au sens d'\italique{ouvrir une ceinture}\nc{Voir LSJ, I. b. Attesté par Héliodore le médecin (I\expo{er} ou \nieme{II} s. ap. J.-C.). De longs fragments nous en ont été conservés par Oribase (\italique{Collectiones medicae} XLVIII, 58.5) : cf. les \italique{Œuvres d'Oribase}, t. 4, p. 320 (1862), de Bussemaker et Daremberg. À propos d'un patient à qui l'on applique un bandage, appelé ceinture : τῇ ζώνῃ \textbf{περιχάσκων} ζώννυται ὁ πάσχων, καὶ τότε ὁ εὐθὺς τελαμὼν ἀπὸ τῆς ὀσφύος ἄγεται ἐπὶ τὴν ἕδραν καὶ τὸν περίναιον, « On entoure le malade avec la ceinture, en l'\italique{ouvrant} largement ; puis on amène la bande droite des lombes au siège et au périnée » (trad. Daremberg).}.\\
Il est remarquable de noter que la tradition millénaire du texte d'Aristophane ait conservé la leçon ἐχύμησαν avec un η jusqu'au scribe du \italique{Ravennas} ; et ce en dépit de manuscrits anciens fautifs portant ἐχύμισαν et de commentaires afférents (ἔγχυμον ἐποίησαν) -- ayant pénétré jusqu'aux marges du \italique{Ravennas} -- suggérant la présence d'un verbe *χυμίζω. Preuve encore une fois que l'œuvre conservée d'Aristophane est, comme l'avait déjà noté Willems, l'une des plus fidèlement transmises avec l'\italique{Iliade} et l'\italique{Odyssée}\nc{Voir l'introduction à ses premières notes sur les \italique{Cavaliers}. L'une des raisons en est que le Comique, au contraire des Tragiques, ne fut que peu rejoué. Les premières éditions critiques de ses œuvres furent sans doute beaucoup plus aisées à constituer que celles d'un Euripide, dont beaucoup plus de versions différentes des pièces circulaient. Victor Coulon, l'éditeur aux Belles Lettres d'Aristophane, abonda dans le sens de Willems : à l’opposé du courant majoritaire de son temps, il s'est astreint à conserver le plus possible les leçons des manuscrits contre les conjectures des philologues. Ce « conservatisme », qu'il défend pour le cas d'Aristophane dans son \italique{Essai sur la méthode de la critique conjecturale appliquée au texte d’Aristophane} (1933) lui valut d’être encore considéré à la fin des années 90 par Sommerstein comme le meilleur éditeur d’Aristophane.}.}
